\documentclass[12pt,a4paper]{article}
\usepackage[margin=2.5cm]{geometry}
\usepackage{graphicx}
\usepackage[colorlinks=true,linkcolor=blue,urlcolor=blue]{hyperref}
\usepackage{listings}
\usepackage{xcolor}
\usepackage{amsmath}
\usepackage{float}
\usepackage{caption}
\usepackage{booktabs}
\usepackage{enumitem}
\usepackage{parskip}
\usepackage{titlesec}
\usepackage{fancyhdr}
\usepackage{subcaption}
\definecolor{codebg}{RGB}{245,245,245}
\definecolor{codeframe}{RGB}{200,200,200}
\definecolor{gray}{RGB}{100,100,100}
\lstdefinestyle{bash}{basicstyle=\ttfamily\small,backgroundcolor=\color{codebg},frame=single,rulecolor=\color{codeframe},breaklines=true,columns=fullflexible,keepspaces=true,showstringspaces=false}
\lstdefinestyle{json}{basicstyle=\ttfamily\footnotesize,backgroundcolor=\color{codebg},frame=single,rulecolor=\color{codeframe},breaklines=true,columns=fullflexible,keepspaces=true,showstringspaces=false,numbers=left,numberstyle=\tiny\color{gray},numbersep=5pt}
\hypersetup{pdftitle={Village Pond Planning -- Phase 1},pdfauthor={Snehal Suhane}}
\pagestyle{fancy}
\fancyhf{}
\lhead{Village Pond Planning System}\rhead{Phase 1 Report}\rfoot{\thepage}
\setlist[itemize]{noitemsep,topsep=4pt}
\setlist[enumerate]{noitemsep,topsep=4pt}
\begin{document}
\begin{titlepage}
\centering\vspace*{2cm}
{\Huge\bfseries Village Pond Planning System\par}\vspace{0.5cm}
{\Large Phase 2: Pond Catchment Analysis Backend\par}\vspace{2cm}
{\large\bfseries Snehal Suhane\par}\vspace{0.3cm}
{\large\bfseries 12342100\par}\vspace{0.3cm}
\hrule\vspace{0.4cm}
\begin{tabular}{ll}
\textbf{GitHub:}&\url{https://github.com/snehalsuhane/pond-planning}\\[4pt]
\textbf{API Base URL:}&\texttt{http://10.1.75.51:5273}\\[4pt]
\textbf{Endpoint:}&\texttt{POST /api/analyzeContour}\\
\end{tabular}
\vspace{0.4cm}\hrule
\end{titlepage}
\tableofcontents\newpage

\section{Introduction}

Rainwater harvesting through village farm ponds is one of the most cost-effective interventions for improving agricultural water security in rain-fed regions of India. Selecting a pond site manually from a contour map is time-consuming, subjective, and difficult to replicate. This project automates that process: given any contour map in KML or KMZ format, the system identifies the best pond locations and quantifies the catchment area that will supply each pond.

The implementation is a \textbf{Flask REST API} backed by a modular Python analysis pipeline. Each module can be independently tested, configured, and extended. The key design principle throughout is \emph{generalisation}: every algorithm is parameterised by the input data, not by assumptions about the specific survey area.

\section{System Architecture \& Design Decisions}

\subsection{Processing Pipeline}

The full pipeline is stateless and driven entirely by the uploaded file:

\begin{enumerate}
\item \textbf{Parse}: KML/KMZ $\rightarrow$ list of \texttt{\{elevation, coordinates\}} dicts
\item \textbf{Terrain validation}: elevation range, point density, geographic bounds
\item \textbf{UTM projection}: WGS84 $\rightarrow$ metres (auto-zone selection)
\item \textbf{DEM interpolation}: scattered contour points $\rightarrow$ regular elevation grid
\item \textbf{Slope computation}: gradient magnitude of the DEM
\item \textbf{Hydrology}: D8 flow direction $\rightarrow$ accumulation $\rightarrow$ channel mask
\item \textbf{Pond candidate ranking}: TPI multi-criteria scoring, greedy selection
\item \textbf{Catchment delineation}: BFS upstream trace on the D8 grid
\item \textbf{Response}: structured JSON
\end{enumerate}

\subsection{Why KML/KMZ? Custom Parser vs.\ Third-Party Libraries}

The system accepts both raw KML (XML) and KMZ (zipped KML). Rather than depending on heavy geo-libraries such as \texttt{geopandas} or \texttt{GDAL}, I wrote a lightweight XML parser using Python's built-in \texttt{xml.etree.ElementTree}. This has several advantages:

\begin{itemize}
\item \textbf{Fewer dependencies}: no system-level GDAL installation required
\item \textbf{Namespace-aware}: handles both the standard \texttt{opengis.net/kml/2.2} namespace and non-namespaced KML files
\item \textbf{KMZ support}: transparently extracts the inner \texttt{.kml} from a ZIP archive
\item \textbf{Explicit error messages}: custom \texttt{KMLParseError} gives precise feedback on malformed inputs
\end{itemize}


\subsection{UTM Projection: Why Not Geographic Coordinates?}

All analysis is performed in a projected coordinate system (Universal Transverse Mercator, UTM) rather than geographic degrees. The reasons are fundamental:

\begin{itemize}
\item \textbf{Metric units required}: DEM grid spacing, slope angles, catchment areas, and minimum candidate separation distances are all expressed in metres. Computing these in degrees is geometrically incorrect near non-equatorial latitudes.
\item \textbf{Isotropic grid cells}: UTM produces square cells; degree-based grids have rectangular cells (longitude degrees are shorter than latitude degrees), distorting slope and flow direction calculations.
\item \textbf{Auto-zone selection}: \texttt{utils/projection.py} determines the correct UTM zone from the bounding box of the input data using \texttt{pyproj}, so no manual configuration is needed for any region of the world.
\end{itemize}

\subsection{DEM Interpolation: Radial Basis Function}

Contour lines provide elevation values along irregular curves, not on a regular grid. To compute slope and flow direction we need a \textbf{Digital Elevation Model (DEM)}: a regular grid of elevation values. I use \textbf{scipy's Radial Basis Function (RBF) interpolation} for this step.

\textbf{Why RBF over alternatives?}

\begin{itemize}
\item \textbf{vs.\ Linear/bilinear interpolation}: Linear methods produce flat ``facets'' between contour lines and sharp ridges at contours; RBF produces smooth, natural-looking terrain that better represents the ground surface.
\item \textbf{vs.\ Inverse Distance Weighting (IDW)}: IDW exhibits the ``bull's-eye'' artefact around data points; RBF minimises this by fitting a globally smooth surface.
\item \textbf{vs.\ Kriging}: Kriging is statistically optimal but requires variogram fitting and is computationally expensive for 200k+ input points. Scipy RBF achieves comparable smoothness with far less complexity.
\item \textbf{Adaptive resolution}: The DEM cell size is set proportionally to the average inter-contour spacing, so fine-grained (1\,m interval) and coarse (5\,m interval) maps are both handled correctly.
\end{itemize}

\subsection{D8 Flow Direction: Why Deterministic Single-Flow?}

Flow direction is the foundation of the catchment delineation algorithm. I am using the \textbf{D8 (Deterministic 8-direction)} algorithm, which directs each cell to its single steepest downhill neighbour among the 8 cardinal/diagonal directions.

\textbf{Why D8 over alternatives?}

\begin{itemize}
\item \textbf{vs.\ D-Infinity (Dinfinity)}: D-Infinity distributes flow proportionally between two downhill neighbours, producing fractional flow weights. This is more realistic for dispersed overland flow but complicates the BFS catchment trace significantly. For pond planning, we want a \emph{definitive} upstream contributing area, which D8 provides directly.
\item \textbf{vs.\ Multiple Flow Direction (MFD)}: MFD spreads flow across all downhill neighbours. Again more realistic for surface runoff modelling, but yields a probabilistic catchment that is harder to interpret as a single bounded polygon.
\item \textbf{Reproducibility}: D8 is deterministic: the same DEM always produces the same flow direction grid. This is important for consistent API responses and testability.
\item \textbf{Established standard}: D8 uses ArcGIS-standard direction codes (E=1, SE=2, S=4, SW=8, W=16, NW=32, N=64, NE=128), making results interpretable with standard GIS tools.
\end{itemize}

\subsection{Pond Candidate Identification: TPI Multi-Criteria Scoring}

\subsubsection{Topographic Position Index (TPI)}

The \textbf{Topographic Position Index} is defined as:
\begin{equation}
    \text{TPI}(c) = z(c) - \frac{1}{|W|}\sum_{n \in W} z(n)
\end{equation}
where $z(c)$ is the elevation of cell $c$ and $W$ is a circular neighbourhood window of radius $r_w$. A \emph{negative} TPI indicates the cell is lower than its surroundings: a depression or valley bottom, which is ideal for water harvesting.

\textbf{Why TPI over raw elevation?} A cell at low absolute elevation may be on a gentle hillside with no local convergence. TPI distinguishes true depressions from low-lying flat areas, making it a far better predictor of natural water collection than elevation alone.

\subsubsection{Multi-Criteria Score}

The composite score combines three normalised components:
\begin{equation}
    \text{score}(c) = w_e \cdot \tilde{S}_e(c) + w_s \cdot \tilde{S}_s(c) + w_d \cdot \tilde{S}_d(c)
\end{equation}

\begin{itemize}
\item $\tilde{S}_e$: \textbf{Elevation score}: lower-elevation sites have more terrain above them to collect runoff. Cells above a slope threshold ($>8°$) are masked out before scoring.
\item $\tilde{S}_s$: \textbf{Slope score}: shallower slopes reduce excavation cost and embankment requirements. Scored as $1 - \text{slope}/\text{max\_slope}$.
\item $\tilde{S}_d$: \textbf{Depression score}: derived from TPI; more negative TPI = stronger depression = higher score. Cells above the slope threshold receive $S_d = 0$.
\item Default weights: $w_e=0.30$, $w_s=0.40$, $w_d=0.30$. All weights are API-configurable.
\end{itemize}

All three components are normalised to $[0,1]$ before weighting, so the score is independent of the absolute scale of the input DEM.

\subsubsection{Greedy Candidate Selection with Minimum Separation}

After scoring all cells, candidates are selected greedily: the highest-scoring cell is selected, then all cells within a minimum Euclidean distance (default 100\,m) are masked, and the process repeats until $N$ candidates are found.

\textbf{Why greedy over clustering?}
\begin{itemize}
\item Clustering (k-means, DBSCAN) requires tuning cluster parameters and may merge two genuinely distinct good sites into one cluster.
\item The greedy approach is $\mathcal{O}(N \cdot k)$ where $k$ is the number of candidates: efficient even on large grids.
\item The minimum separation distance is a physically meaningful parameter (minimum distance between two ponds) rather than an abstract cluster bandwidth.
\end{itemize}

\subsection{Catchment Delineation: BFS on the D8 Grid}

Given a pour point $(r_0, c_0)$, the catchment is all cells that ultimately drain into it under the D8 model. We find these with a \textbf{Breadth-First Search (BFS)} over the reverse D8 graph:

\begin{lstlisting}[style=bash]
# Build reverse lookup: for each direction code, which offset points TO this cell?
REVERSE_D8 = {(dr, dc): reverse_code, ...}

queue = deque([(r0, c0)])
mask[r0, c0] = True
while queue:
    r, c = queue.popleft()
    for (dr, dc), rev_code in REVERSE_D8.items():
        nr, nc = r + dr, c + dc
        if in_bounds(nr, nc) and not mask[nr, nc]:
            if flow_direction[nr, nc] == rev_code:
                mask[nr, nc] = True
                queue.append((nr, nc))
\end{lstlisting}

The catchment area is:
\begin{equation}
    A = N_{\text{cells}} \times r^2 \quad [\text{m}^2]
\end{equation}
where $r$ is the DEM resolution. The polygon boundary is extracted from the raster mask using \texttt{contourpy} (marching squares) and re-projected from UTM to WGS84.

\textbf{Why BFS over watershed segmentation tools?}
\begin{itemize}
\item No additional dependencies (no GRASS, SAGA, or TauDEM required)
\item The algorithm is transparent, testable, and directly tied to the D8 grid produced in the hydrology step
\item BFS is $\mathcal{O}(\text{rows} \times \text{cols})$: linear in the grid size, efficient for grids up to thousands of cells on a side
\end{itemize}

\section{API Documentation}

\subsection{Endpoint}

\begin{center}
\large\textbf{\texttt{POST /api/analyzeContour}}
\end{center}

\begin{table}[H]
\centering
\begin{tabular}{@{}lll@{}}
\toprule
\textbf{Parameter} & \textbf{Type} & \textbf{Description} \\
\midrule
\texttt{contour\_map} & \texttt{multipart/form-data} & KML or KMZ contour map \\
\bottomrule
\end{tabular}
\caption{Request parameters}
\end{table}

\subsection{Example Request}

\begin{lstlisting}[style=bash]
curl -X POST http://10.1.75.51:5273/api/analyzeContour \
     -F "contour_map=@contours_1m.kml"
\end{lstlisting}

\subsection{Response Schema}

\begin{lstlisting}[style=json, caption={Annotated JSON response (abbreviated)}]
{
  "status": "success",
  "filename": "contours_1m.kml",
  "terrain": {
    "contour_count": 1355,        // Number of contour lines parsed
    "min_elevation_m": 267.0,
    "max_elevation_m": 298.0,
    "contour_interval_m": 1.0,    // Dominant vertical interval
    "contour_interval_uniform": true,
    "total_points": 285619,
    "bounds": {
      "min_lon": 81.2799, "max_lon": 81.3117,
      "min_lat": 21.2369, "max_lat": 21.2769
    },
    "crs": { "epsg": 32644, "name": "WGS 84 / UTM zone 44N" }
  },
  "dem": {
    "resolution_m": 6.0,          // Grid cell size in metres
    "shape": [439, 542],          // [rows, cols]
    "nan_fraction": 0.0,          // 0 = fully interpolated
    "elevation_min": 267.0, "elevation_max": 298.0,
    "bounds": { "min_x": ..., "max_x": ..., "min_y": ..., "max_y": ... },
    "slope": {
      "slope_min_deg": 0.0, "slope_max_deg": 51.3, "slope_mean_deg": 3.4
    }
  },
  "pond_candidates": [
    {
      "rank": 1,                  // 1 = best site
      "latitude": 21.259564,
      "longitude": 81.300134,
      "elevation_m": 274.1,
      "slope_deg": 4.3,
      "tpi": -8.26,              // Negative = depression
      "score": 0.102,            // Lower = better rank
      "criteria": {
        "elevation_score": 0.031,
        "slope_score": 0.036,
        "depression_score": 0.035
      },
      "catchment": {
        "area_m2": 33264.0,
        "area_ha": 3.33,
        "area_km2": 0.033264,
        "polygon": [[81.2989, 21.2571], ...]  // WGS84 [lon, lat] ring
      }
    }
    // ... up to 10 candidates, ranked by score ascending
  ],
  "hydrology": {
    "noflow_count": 35753,        // D8 pit/flat cells (code=0)
    "acc_max": 924,               // Max flow accumulation (cells)
    "acc_mean": 7.3,
    "channel_threshold": 81.6,    // Accumulation threshold for channels
    "channel_cell_count": 2380,
    "channel_fraction": 0.01      // Fraction of grid that is channel
  }
}
\end{lstlisting}

\subsection{Error Responses}

\begin{table}[H]
\centering
\begin{tabular}{@{}cll@{}}
\toprule
\textbf{HTTP} & \textbf{Condition} & \textbf{Description} \\
\midrule
400 & Missing \texttt{contour\_map} field & No contour\_map key in multipart body \\
400 & Empty filename & File part has no filename \\
415 & Wrong extension & Not \texttt{.kml} or \texttt{.kmz} \\
422 & Validation failure & Malformed KML or terrain validation fails \\
500 & Internal error & Unexpected server-side exception \\
\bottomrule
\end{tabular}
\caption{Error response HTTP codes}
\end{table}

All errors return: \texttt{\{ "status": "error", "error": "<message>" \}}


\section{Demonstration: Sample Contour Map}

\subsection{Input Data}

The provided sample file \texttt{contours\_1m.kml} covers a rural area near Bilaspur, Chhattisgarh, India (approx.\ $21.24^\circ$N--$21.28^\circ$N, $81.28^\circ$E--$81.31^\circ$E). It contains \textbf{1355 contour lines} at a uniform 1\,m vertical interval spanning 267\,m to 298\,m elevation (31\,m total relief), with 285,619 coordinate points. The API accepts this file unmodified with no configuration changes.

\subsection{Terrain and DEM}

\begin{figure}[H]
    \centering
    \includegraphics[width=\textwidth]{report_dem.png}
    \caption{\textbf{Terrain visualisation for the sample contour map.}
    \textit{Left:} Hillshaded elevation surface (315$^\circ$ azimuth, 45$^\circ$ altitude light source): hill shading makes valley and ridge structure immediately apparent.
    \textit{Centre:} Slope map in degrees (yellow = flat $<2^\circ$, red = steep $>15^\circ$): the flat valley bottoms are the candidate pond zones.
    \textit{Right:} Contour lines reconstructed from the interpolated DEM, confirming the interpolation accurately reproduces the input survey data.}
    \label{fig:dem}
\end{figure}

The DEM is a $439 \times 542$ grid at 6\,m/cell. Mean slope is 3.4$^\circ$; steep gully walls reach 51.3$^\circ$. The flat valley bottoms visible in the slope panel are exactly where the pond scoring algorithm concentrates candidates.

\subsection{Pond Candidate Results}

\begin{figure}[H]
    \centering
    \includegraphics[width=\textwidth]{report_pond.png}
    \caption{\textbf{Top 10 pond candidates.}
    \textit{Left:} Hillshaded DEM with rainbow-ranked sites (red star = \#1 best). Score and rank are annotated at each site.
    \textit{Centre:} Slope map: all candidates sit in low-slope zones ($<5^\circ$), confirming the scoring correctly penalises steep cells.
    \textit{Right:} TPI zoom ($\pm$400\,m) around the \#1 site; deep blue = strong depression, confirming the site is a genuine topographic low.}
    \label{fig:candidates}
\end{figure}

\begin{table}[H]
\centering
\small
\begin{tabular}{@{}rrrrrrl@{}}
\toprule
\textbf{Rank} & \textbf{Latitude} & \textbf{Longitude} & \textbf{Elev (m)} & \textbf{Slope ($^\circ$)} & \textbf{TPI} & \textbf{Score} \\
\midrule
1  & 21.259564 & 81.300134 & 274.1 & 4.3 & $-$8.26 & 0.102 \\
2  & 21.252756 & 81.287167 & 272.0 & 0.0 & $-$5.03 & 0.166 \\
3  & 21.255305 & 81.286304 & 270.0 & 0.1 & $-$4.46 & 0.168 \\
4  & 21.251710 & 81.296475 & 270.0 & 0.0 & $-$3.66 & 0.196 \\
5  & 21.258070 & 81.286136 & 269.0 & 0.8 & $-$3.47 & 0.200 \\
6  & 21.249769 & 81.289995 & 267.0 & 0.8 & $-$2.74 & 0.206 \\
7  & 21.252591 & 81.288497 & 269.0 & 0.0 & $-$3.02 & 0.210 \\
8  & 21.263226 & 81.282965 & 267.0 & 0.0 & $-$2.27 & 0.217 \\
9  & 21.258696 & 81.300480 & 279.0 & 0.0 & $-$5.46 & 0.218 \\
10 & 21.244784 & 81.288828 & 268.0 & 0.0 & $-$2.42 & 0.222 \\
\bottomrule
\end{tabular}
\caption{Top 10 pond candidates. All TPI values are negative, confirming every site is a topographic depression. Score is lower-is-better (0 = perfect match on all criteria).}
\label{tab:candidates}
\end{table}

\subsection{Catchment Delineation Results}

\begin{figure}[H]
    \centering
    \includegraphics[width=\textwidth]{report_catchment.png}
    \caption{\textbf{Catchment delineation for all 10 candidates.}
    Each coloured zone shows the upstream area that drains into that pour point under the D8 model.
    Candidates with small true catchments ($<5000$\,m$^2$) are displayed with a minimum-radius expanded zone for visibility; the legend shows the actual computed area in m$^2$.
    The \#1 candidate (red star, top-right cluster) dominates with a 33,264\,m$^2$ catchment, more than 20$\times$ larger than any other site.}
    \label{fig:catchment}
\end{figure}

\begin{table}[H]
\centering
\small
\begin{tabular}{@{}rrrl@{}}
\toprule
\textbf{Rank} & \textbf{Catchment (m$^2$)} & \textbf{Catchment (ha)} & \textbf{Notes} \\
\midrule
1  & 33\,264 & 3.33 & Strong D8 upstream convergence into valley depression \\
2  & 36      & 0.00 & Single no-flow pit cell; no upstream inflow \\
3  & 216     & 0.02 & Small 6-cell sub-basin \\
4  & 36      & 0.00 & No-flow pit cell \\
5  & 36      & 0.00 & No-flow pit cell \\
6  & 1\,512  & 0.15 & Moderate upstream convergence \\
7  & 1\,656  & 0.17 & Moderate upstream convergence \\
8  & 36      & 0.00 & No-flow pit cell \\
9  & 36      & 0.00 & No-flow pit cell \\
10 & 36      & 0.00 & No-flow pit cell \\
\bottomrule
\end{tabular}
\caption{Catchment areas for the sample map. The 36\,m$^2$ entries (1 grid cell at 6\,m resolution) are D8 pit cells where no surrounding cell drains inward.}
\label{tab:catchment}
\end{table}

The \textbf{\#1 candidate} (lat 21.2596, lon 81.3001, elev 274.1\,m, slope 4.3$^\circ$, TPI $-$8.26) has the deepest depression score and the largest verified catchment of \textbf{33,264\,m$^2$ (3.33\,ha)}. This is the recommended site for a village farm pond on this terrain: it sits in the strongest natural depression, has a buildable slope, and drains the largest upstream area.

\section{Generalisation and Extensibility}

A core requirement is that the system must generalise to other contour maps without code changes. The following design decisions ensure this:

\begin{itemize}
\item \textbf{No hard-coded coordinates or elevations}: every threshold, grid size, and score is derived from the statistics of the uploaded file.
\item \textbf{Automatic UTM zone selection}: works anywhere on Earth; tested with maps from central India and designed to handle any KML input from other regions.
\item \textbf{Adaptive DEM resolution}: the grid spacing is computed as a fraction of the average inter-contour distance, so both 1\,m and 5\,m contour intervals produce appropriately-resolved DEMs.
\item \textbf{Configurable scoring weights and candidate count}: API callers can override defaults via request parameters, enabling per-region tuning without changing server code.
\item \textbf{Modular pipeline}: each stage (\texttt{terrain}, \texttt{dem}, \texttt{hydrology}, \texttt{pond}, \texttt{catchment}) is an independent module with its own test suite. Future phases can replace or augment individual stages (e.g., add soil permeability, seasonal runoff coefficients, or road accessibility layers) without touching the rest of the pipeline.
\end{itemize}

\section{Test Coverage}

\begin{table}[H]
\centering
\begin{tabular}{@{}llp{8cm}@{}}
\toprule
\textbf{Module} & \textbf{Tests} & \textbf{Covers} \\
\midrule
\texttt{test\_contour\_route.py} & 5  & HTTP layer; full real pipeline (no mocks) \\
\texttt{test\_kml\_parser.py}    & 20 & KML/KMZ parsing, namespaces, edge cases \\
\texttt{test\_terrain.py}        & 27 & Stats, contour intervals, all validation errors \\
\texttt{test\_projection.py}     & 33 & UTM zone selection, coordinate projection, pipeline \\
\texttt{test\_dem.py}            & 28 & DEM shape, range, NaN fraction, cache reuse \\
\texttt{test\_pond.py}           & 26 & Scoring normalisation, top-N, TPI depression, criteria sum \\
\texttt{test\_hydrology.py}      & 32 & D8 direction codes, accumulation, channel detection \\
\texttt{test\_catchment.py}      & 7  & BFS upstream trace, area units, polygon WGS84 bounds \\
\midrule
\textbf{Total} & \textbf{178} & \textbf{All passing} \\
\bottomrule
\end{tabular}
\caption{Test suite summary. Run with: \texttt{PYTHONPATH=. pytest tests/ -v}}
\end{table}

\section{Running the API Locally}

\begin{lstlisting}[style=bash]
# 1. Clone
git clone https://github.com/snehalsuhane/pond-planning.git
cd pond-planning

# 2. Virtual environment
python3 -m venv .venv && source .venv/bin/activate
pip install -r requirements.txt

# 3. Start server
PYTHONPATH=. python app.py
#  * Running on http://10.1.75.51:5273

# 4. Run test suite
PYTHONPATH=. pytest tests/ -v   # 178 tests, all passing

# 5. Upload sample map
curl -X POST http://10.1.75.51:5273/api/analyzeContour \
     -F "contour_map=@contours_1m.kml"
\end{lstlisting}

\section{Live API Demonstration}

The API is deployed and accessible at:

\begin{center}
    \large\texttt{http://10.1.75.51:5273/api/analyzeContour}
\end{center}

\subsection{curl Request from Local Machine}

The following command was used to submit the sample contour map to the deployed server from a local machine on the same network:

\begin{lstlisting}[style=bash]
curl -X POST http://10.1.75.51:5273/api/analyzeContour \
     -F "contour_map=@contours_1m.kml"
\end{lstlisting}

\subsection{API Response (Screenshot)}

Figure~\ref{fig:curl} shows a screenshot of the terminal output after sending the request. The full JSON response is truncated in the screenshot as it is several hundred lines long; the visible portion shows the \texttt{status}, \texttt{filename}, \texttt{terrain}, and the beginning of the \texttt{dem} block, confirming the server received the file, parsed 1355 contours, and began returning structured results.

\begin{figure}[H]
    \centering
    % ---------------------------------------------------------------
    % Place your curl screenshot here:
    % \includegraphics[width=\textwidth]{curl_response.png}
    % ---------------------------------------------------------------
    \fbox{\parbox[c][9cm][c]{0.95\textwidth}{\centering
        \textit{[Insert screenshot of curl response here]}\\[0.5cm]
        \small Filename: \texttt{curl\_response.png}
    }}
    \caption{\textbf{curl response from the deployed API at \texttt{http://10.1.75.51:5273}.}
    The terminal shows the beginning of the JSON response from \texttt{POST /api/analyzeContour}
    with the sample contour map \texttt{contours\_1m.kml}.
    The full response (top 10 candidates with catchment polygons) is several hundred lines;
    only the initial portion is visible in the screenshot.}
    \label{fig:curl}
\end{figure}

\section{Conclusion}

This phase delivers a fully generalised REST API for village pond site selection and catchment estimation. The key contributions are:

\begin{enumerate}
\item A robust KML/KMZ ingest pipeline with strict validation and informative error messages
\item An RBF-based DEM interpolation that produces smooth, accurate elevation grids from irregular contour data
\item A D8 hydrological model that preserves natural depressions: the very features that make good pond sites: rather than eliminating them via pit-filling
\item A TPI-based multi-criteria scoring system that combines topographic position, slope suitability, and relative elevation into a single, interpretable score
\item A BFS catchment delineation that computes exact upstream contributing areas with no external GIS dependencies
\item A clean, well-documented JSON API with 178 automated tests verifying correctness at every stage
\end{enumerate}

The system is production-ready for the sample map and designed to handle arbitrary KML/KMZ contour inputs in future phases without modification.

\end{document}
